% Source pins (SHA-256): PifoStatement.lean=fe9475ebfa2c462d27cb0a99781926074e1281396f2c9e8d32d6dfca788d814e; pifo-elegant-v4.md=356ea38146ced437e20200db04e2ce220c73f9ac9fa892b746f783b8a2d49ab8; pifo-nbe.md=5f7efd28922efc5dedad29a375f030ad6e3fe29f96ba94e77ec308f521541f14; pifo-normalization.md=2d287de7924aa30479218da3e0a9a31577ecb97c5a34395c089f7abed98f61e7; pifo-eta-decision.md=50159fe6aaf8273fc1d811f941de4e2246f7ec84e2270fefc1144c0a890e2ec9; pifo-confluence.md=58274d7fa628cbd04ae265a415cd11bce671e950d53a1653c6e23fe5ee1342ec; pifo-noinduction-codex2.md=97eef33723fdd15eff0003d02ce1cd387a1b1dcd86a42f89302421f0fa614933; pifo-verify-all/VERIFICATION.md=c97ed5daa11c97f71a5d0e3867f96ddabe74438f1f598c049050c735a96b1887.
\section{Introduction}
\label{sec:introduction}

A scheduler can be questioned in two ways.  A \emph{flush} experiment first
pushes a finite word of packets and then pops until every pushed packet has
been returned.  An \emph{interleaved} experiment may alternate pushes and pops
arbitrarily.  Interleaving is apparently more informative: an early pop can
change which packets remain together, and later arrivals then meet a residual
state that no batch drain exposes directly.  This raises a basic extensional
question for packet schedulers:

\begin{quote}
If two schedulers produce the same drain for every input word, must they
produce the same observations for every interleaving of pushes and pops?
\end{quote}

For arbitrary stateful assignment policies, the answer is no.  For stateless
PIFO trees, the answer is yes.  Here a PIFO is a priority queue ordered by rank,
with FIFO order among equal ranks.  A PIFO tree places such a queue at every
node: an internal queue chooses a child, and a leaf queue chooses a packet.  A
push follows a root-to-leaf path and inserts one ranked entry at every visited
node.  \emph{Stateless} means that the packet type alone fixes this entire
annotated path.  It may not depend on the current queues, the arrival number,
or earlier operations.

Our main theorem says that, over every finite packet alphabet and arbitrary
finite tree topologies, flush equivalence implies interleaved equivalence for
valid stateless assignments.  Since a flush is itself an operation word, the
reverse implication is immediate.  Thus flushes and general interleavings
coincide extensionally on this class.

This result gives precise content to the practice of testing a scheduler by
batch drains.  As stated, its premise quantifies over all finite input words:
the implication first says that no additional observational principle is
gained by admitting residual-state experiments once all batch behavior is
known.  The later canonical analysis strengthens this with a finite test
theorem.  This distinction matters both conceptually and formally.  The main
proof reduces an operational equivalence over mixed commands to equality of a
pure word transformer, without assuming that either scheduler's hidden
topology or root partition can be reconstructed from observations.

The proof has four ingredients.

\begin{enumerate}
\item On two packet types, a stateless PIFO tree is a single stable comparison.
      Two orderings of the binary flush distinguish strict priority in either
      direction from a FIFO tie.
\item A colored-PIFO coupling shows that changing ranks within colors preserves
      the complete online color trace whenever all cross-color comparisons are
      unchanged.
\item The top partitions of two flush-equivalent trees induce an incidence
      grid.  Binary comparisons agree wherever both partitions separate a
      pair.  A three-cell diamond rules out local comparison cycles, and a
      fixed-head-chord induction rules out cycles of every length.  The result
      is one total preorder compatible with both roots across their blocks.
\item Both roots are first swapped to this common preorder.  Each proper block
      is then refined into meet cells and discharged by strong induction on the
      alphabet.  The remaining two selector levels use identical keys and
      collapse occurrence-for-occurrence.
\end{enumerate}

The same local structure yields more than the implication.  Effective
comparisons on pairs and rooted obstruction bits on triples generate a
canonical component-decomposition tree, which is a complete invariant for
interleaved behavior.  On a unit-cost RAM, they give an \(O(k^3h)\) structural
decision procedure, an injective test suite using at most three pushes, and a matching
\(\Theta(k^3)\) lower bound on triple supports among suites whose tests use at
most three types.  An independent eta-style comparator decides equivalence
directly through pair-dependent meet obligations, without constructing either
normal form; it returns a two- or three-push counterexample or a finite
equivalence proof DAG.  The natural active rewrite equations do not provide a
rewrite normal form: collapse/regrouping and leaf explosion create cycles, and
even a canonical tree is reducible.  Nevertheless a compositional bottom-up
fold over restriction-ready pair/triple summaries computes exactly the
canonical form.  A normalization-by-evaluation argument instead evaluates to
the literal flush transformer, reifies its finest factor hierarchy, and gives
a second semantic proof, not a decision algorithm.  Finally, a separate
meet-promotion argument removes induction on the number of packet types; by
finite support, the observational theorem applies to arbitrary type alphabets
over finite topologies, without asserting a global infinite-alphabet normal
form.

Statelessness is sharp as a hypothesis on the class.  We exhibit two flat,
arrival-dependent assignment policies whose flush transformer is identical on
every word but whose outputs differ after
\(\push;\popop;\push;\push;\popop\).  The example isolates the failure mode:
an internal reference selects a leaf, not the packet occurrence that created
that reference.  A prior pop can remove the packet that would otherwise have
been ``hijacked'' by a later reference.  A fixed type-to-path assignment rules
out exactly the overtaking freedom used by this construction.

Besides proving the equivalence theorem, the argument gives a useful semantic
picture of stateless PIFO trees.  They are arbitration hierarchies.  Each pair
of types meets at a first decisive queue, while the higher queues schedule only
coarser child colors.  Unary forwarding depth is operationally invisible, and
consecutive selectors with the same preorder can be collapsed.  This picture
is developed in \autoref{sec:discussion} after the formal proof.

The mathematical model in \autoref{sec:model} is a direct presentation of the
frozen Lean statement.  \autoref{sec:main-proof} gives the complete proof,
including the binary recovery table, colored coupling, meet construction, and
collapse.  \autoref{sec:stateful} works through the stateful separation.
The later mechanization section records statement hashes, clean-slate proof
checking, per-artifact independence, axiom audits, and the exact scope of the
kernel witnesses.  The
subsequent sections develop canonical normalization, structural and eta-style
decision, the negative rewrite result and positive bottom-up normalizer,
normalization by evaluation, and topology-recursive proofs with their
finite-support extension.
